\section{仿射簇上的多项式函数}\label{sec:polynomial-functions-on-affine-varieties}

第~\ref{sec:vanishing-set-operator}~节说明了每个多项式 $f\in K[x_{1},\ldots,x_{n}]$ 都可用来定义函数 $\mathbb{A}^{n}\to K$, 把 $(a_{1},\ldots,a_{n})$ 映到 $f(a_{1},\ldots,a_{n})$. 若 $X\subseteq\mathbb{A}^{n}$ 是一个仿射簇, 就可以把这样的函数的定义域限制到 $X$, 得到一个新函数
\[
  \begin{aligned}
    f|_{X}:X&\to K\\
    (a_{1},\ldots,a_{n})&\mapsto f(a_{1},\ldots,a_{n}).
  \end{aligned}
\]
由这种方式得到的函数 $F:X\to K$ 称为多项式函数.

\begin{Definition}{多项式函数}\leavevmode\label{def:polynomial-function-on-affine-variety}
  设 $X\subseteq\mathbb{A}^{n}$ 是仿射簇. $X$ 上的\textbf{多项式函数}是满足 $F=f|_{X}$ 的函数 $F:X\to K$, 其中 $f\in K[x_{1},\ldots,x_{n}]$ 为某个多项式.
\end{Definition}

\begin{Example}{抛物线上的多项式函数}\leavevmode\label{ex:polynomial-functions-on-parabola}
  设 $X=\mathcal{V}(y-x^{2})\subseteq\mathbb{A}^{2}$. 函数
  \[
    \begin{aligned}
      F:X&\to K\\
      (a,b)&\mapsto a+b
    \end{aligned}
  \]
  是多项式函数, 因为
  \[
    F=f|_{X}\qquad\text{其中}\qquad f=x+y\in K[x,y].
  \]
  注意, $f=x+y$ 并不是唯一能给出 $F$ 的多项式. 例如, 由于对所有 $(a,b)\in X$ 都有 $a^{2}=b$, 可知 $F=g|_{X}$, 其中 $g=x+x^{2}$, 同时 $F=h|_{X}$, 其中 $h=x+2y-x^{2}$.
\end{Example}

\begin{Example}{坐标函数}\leavevmode\label{ex:coordinate-functions}
  设 $X\subseteq\mathbb{A}^{n}$ 是仿射簇. 对每个 $i\in\{1,\ldots,n\}$, $X$ 上的第 $i$ 个坐标函数是
  \[
    \begin{aligned}
      C_{i}:X&\to K\\
      (a_{1},\ldots,a_{n})&\mapsto a_{i}.
    \end{aligned}
  \]
  坐标函数是多项式函数, 因为 $C_{i}$ 是多项式 $x_{i}\in K[x_{1},\ldots,x_{n}]$ 所对应函数的限制.
\end{Example}

\begin{Example}{空函数}\leavevmode\label{ex:empty-polynomial-function}
  若 $X=\varnothing\subseteq\mathbb{A}^{n}$, 那么只有一个函数 $F:\varnothing\to K$, 即空函数. 此外, 把定义域限制到空集时, 每个函数 $\mathbb{A}^{n}\to K$ 都给出空函数. 特别地, 这意味着空函数是仿射簇 $\varnothing\subseteq\mathbb{A}^{n}$ 上唯一的多项式函数.
\end{Example}

从以上讨论应当很清楚, 构造 $X\subseteq\mathbb{A}^{n}$ 上的多项式函数, 只须选取任意多项式 $f\in K[x_{1},\ldots,x_{n}]$, 考虑 $\mathbb{A}^{n}$ 上相应的函数, 再把其定义域限制到 $X$. 那么, 为什么定义 $X$ 上的多项式函数时, 我们却似乎采用了相反的方式, 从 $F$ 出发, 再寻找一个限制到 $X$ 后等于 $F$ 的 $f$?

这样做的主要原因是: 给定多项式函数 $F:X\to K$, 可能有许多多项式 $f\in K[x_{1},\ldots,x_{n}]$ 满足 $F=f|_{X}$, 我们并不想把它们看成 $X$ 上不同的多项式函数. 这一点在 Example~\ref{ex:polynomial-functions-on-parabola} 中已经很明显, 其中不同的多项式 $f,g,h\in K[x,y]$ 都定义了同一个函数 $F:X\to K$. 这里重要的不是多项式 $f\in K[x_{1},\ldots,x_{n}]$, 我们要研究的是函数 $F:X\to K$.

不过, 从 $F$ 出发也有缺点. 视 $F$ 的定义如何给出, 它是否为某个多项式所对应函数的限制, 可能并不显然. 下一个例子说明了这个现象, 也提醒我们不要过早判断一个函数是不是多项式函数.

\begin{Example}{一个不显然是多项式函数的函数}\leavevmode\label{ex:apparently-nonpolynomial-hyperbola-function}
  设 $X=\mathcal{V}(xy-1)\subseteq\mathbb{A}^{2}$. 由于任意 $(a,b)\in X$ 都有 $a\ne0$, 可以考虑如下定义的函数:
  \[
    \begin{aligned}
      F:X&\to K\\
      (a,b)&\mapsto\frac{1}{a}.
    \end{aligned}
  \]
  输出 $\frac{1}{a}$ 并不是关于 $a$ 和 $b$ 的多项式, 这也许会让人猜测 $F$ 不是多项式函数. 然而, 对所有 $(a,b)\in X$ 都有 $ab-1=0$, 这意味着 $\frac{1}{a}=b$, 从而 $F=f|_{X}$, 其中 $f=y\in K[x,y]$.

  更一般地, 任意形如
  \[
    F(a,b)=\frac{f(a,b)}{a^{j}b^{k}}
  \]
  的函数 $F:X\to K$, 其中 $f\in K[x,y]$, $j,k\in\mathbb{N}$, 都是 $X$ 上的多项式函数 (Exercise~\ref{ex:hyperbola-polynomial-function-despite-denominators}).
\end{Example}

对 $X$ 上多项式函数的集合, 可以用通常的方式加乘函数, 从而赋予它环结构:
\[
  \begin{aligned}
    (F+G)(a_{1},\ldots,a_{n})&=F(a_{1},\ldots,a_{n})+G(a_{1},\ldots,a_{n}),\\
    (F\cdot G)(a_{1},\ldots,a_{n})&=F(a_{1},\ldots,a_{n})\cdot G(a_{1},\ldots,a_{n}).
  \end{aligned}
\]
因此, 从一个仿射簇 $X$ 出发, 可以得到一个与之关联的环. 这个环在代数几何的研究中居于核心地位.

\begin{Definition}{坐标环}\leavevmode\label{def:coordinate-ring}
  设 $X\subseteq\mathbb{A}^{n}$ 是仿射簇. $X$ 的\textbf{坐标环}记为 $K[X]$, 是 $X$ 上所有多项式函数构成的环.
\end{Definition}

加法单位元 $0\in K[X]$ 是对所有 $a\in X$ 取值 $0\in K$ 的常函数, 乘法单位元是对所有 $a\in X$ 取值 $1\in K$ 的常函数. 这两个函数分别来自多项式 $0,1\in K[x_{1},\ldots,x_{n}]$: $0=0|_{X}$ 且 $1=1|_{X}$.

给定仿射簇 $X$, 能否计算 $K[X]$? 换句话说, 能否把 $K[X]$ 与一个更熟悉的环等同起来? 下面的结果把坐标环表示为一个商环, 朝这个方向迈出一步.

\begin{Proposition}{坐标环的商环表示}\leavevmode\label{prop:coordinate-ring-as-quotient}
  若 $X\subseteq\mathbb{A}^{n}$ 是仿射簇, 则有典范环同构
  \[
    K[X]=\frac{K[x_{1},\ldots,x_{n}]}{\mathcal{I}(X)}.
  \]
\end{Proposition}

\begin{Proof}
  根据第一同构定理, 只须找到一个以 $\mathcal{I}(X)$ 为核的典范满同态 $\phi:K[x_{1},\ldots,x_{n}]\to K[X]$. 定义 $\phi$ 为
  \[
    \phi(f)=f|_{X}.
  \]
  注意 $(f+g)|_{X}=f|_{X}+g|_{X}$ 且 $(f\cdot g)|_{X}=f|_{X}\cdot g|_{X}$, 可知 $\phi$ 是环同态. 根据定义, $X$ 上每个多项式函数都来自某个 $f\in K[x_{1},\ldots,x_{n}]$, 所以 $\phi$ 是满射. 最后, $f\in\ker\phi$ 当且仅当 $f|_{X}=0$, 这等价于 $f\in\mathcal{I}(X)$. 由此可见 $\ker\phi=\mathcal{I}(X)$.
\end{Proof}

\begin{Example}{仿射空间的坐标环}\leavevmode\label{ex:coordinate-ring-of-affine-space}
  由 Proposition~\ref{prop:polynomial-function-isomorphism}, 以及 $K$ 是代数闭域从而是无限域这一事实, 可知 $K[\mathbb{A}^{n}]=K[x_{1},\ldots,x_{n}]$.
\end{Example}

\begin{Example}{抛物线的坐标环}\leavevmode\label{ex:coordinate-ring-of-parabola}
  在 Example~\ref{ex:polynomial-functions-on-parabola} 中, 我们已经看到, 对 $X=\mathcal{V}(y-x^{2})\subseteq\mathbb{A}^{2}$, 三个多项式
  \[
    f=x+y,\quad g=x+x^{2},\quad\text{以及}\quad h=x-x^{2}+2y
  \]
  都给出同一个多项式函数 $F:X\to K$. 这反映了在商环
  \[
    K[X]=\frac{K[x,y]}{\mathcal{I}(X)}=\frac{K[x,y]}{\langle y-x^{2}\rangle}
  \]
  中有 $[f]=[g]=[h]$. 验算如下: $f-g=y-x^{2}\in\langle y-x^{2}\rangle$ 且 $g-h=2x^{2}-2y\in\langle y-x^{2}\rangle$.
\end{Example}

对一个仿射簇 $X\subseteq\mathbb{A}^{n}$, Definition~\ref{def:coordinate-ring} 和 Proposition~\ref{prop:coordinate-ring-as-quotient} 给出了坐标环 $K[X]$ 的两种不同刻画. 记住这两种解释很重要: 坐标环的元素既应看作函数 $F:X\to K$, 也应看作 $K[x_{1},\ldots,x_{n}]$ 中多项式的等价类. 对任意 $f\in K[x_{1},\ldots,x_{n}]$, Proposition~\ref{prop:coordinate-ring-as-quotient} 中的典范同构将多项式函数 $f|_{X}$ 与等价类 $[f]$ 等同起来.

如本章开篇所说, 我们现在已经给出一种从几何走向代数的新方法, 即为一个仿射簇 $X$ 配上它的坐标环 $K[X]$. 按照几何与代数之间的转换应是一本双向词典这一思想, 我们现在要问从仿射簇到其坐标环的转换能否逆转. 更确切地说, 给定一个环 $R$, 是否存在仿射簇 $X$ 使 $R=K[X]$? 我们将看到, 如果 $R$ 是有限生成的既约 $K$-代数, 答案就是肯定的. 接下来三节将定义并研究这些术语.

\BookExerciseHeading
\begin{BookExercises}
  \item\label{ex:coordinate-ring-of-empty-set} 描述空集上的多项式函数环. $0$ 是什么? $1$ 又是什么?
  \item\label{ex:distinct-polynomials-restrict-to-same-function} 设 $X=\mathcal{V}(x^{2}+y^{2}-2z^{2})\subseteq\mathbb{A}^{3}$. 列出 $K[x,y,z]$ 的三个不同元素, 它们限制后给出 $K[X]$ 中同一个多项式函数; 再列出 $K[x,y,z]$ 的两个元素, 它们限制后给出 $K[X]$ 中不同的多项式函数.
  \item\label{ex:hyperbola-polynomial-function-despite-denominators} 设 $X=\mathcal{V}(xy-1)\subseteq\mathbb{A}^{2}$. 证明任意形如
    \[
      F(a,b)=\frac{f(a,b)}{a^{j}b^{k}}
    \]
    的函数 $F:X\to K$, 其中 $f\in K[x,y]$, $j,k\in\mathbb{N}$, 都是多项式函数.
  \item\label{ex:coordinate-ring-domain-iff-irreducible} 设 $X\subseteq\mathbb{A}^{n}$ 是仿射簇.
    \begin{enumerate}[label=(\alph*),ref=(\alph*)]
      \item\label{item:coordinate-ring-domain-iff-irreducible} 证明 $K[X]$ 是整环当且仅当 $X$ 不可约.
      \item 作为 \ref{item:coordinate-ring-domain-iff-irreducible} 的一个例子, 设 $X=\mathcal{V}(x^{2}-xy)\subseteq\mathbb{A}^{2}$. 找到两个仿射簇 $X_{1}\subsetneq X$ 和 $X_{2}\subsetneq X$ 使 $X=X_{1}\cup X_{2}$, 从而证明 $X$ 可约. 然后, 找到 $K[X]$ 中乘积为零的两个非零函数, 验证 $K[X]$ 不是整环.
    \end{enumerate}
  \item\label{ex:coordinate-ring-of-single-point} 设 $X=\{p\}\subseteq\mathbb{A}^{n}$ 是一个单点. 证明由 $\phi(F)=F(p)$ 定义的函数 $\phi:K[X]\to K$ 是环同构.
  \item\label{ex:coordinate-ring-of-finite-points} 设 $p_{1},\ldots,p_{m}\in\mathbb{A}^{n}$ 是 $\mathbb{A}^{n}$ 中互不相同的点, 并令 $X=\{p_{1},\ldots,p_{m}\}$. 证明函数
    \[
      \begin{aligned}
        \phi:K[X]&\to\overbrace{K\oplus\cdots\oplus K}^{m}\\
        \phi(F)&=(F(p_{1}),\ldots,F(p_{m}))
      \end{aligned}
    \]
    是环同构. (回顾一下, 直和的加法和乘法按分量定义.)
  \item\label{ex:equal-coordinate-functions-on-infinite-variety}
    \begin{enumerate}[label=(\alph*),ref=(\alph*)]
      \item\label{item:equal-coordinate-functions-example} 举出一个无限仿射簇 $X\subseteq\mathbb{A}^{3}$, 使三个坐标函数都是同一个多项式函数.
      \item 证明 \ref{item:equal-coordinate-functions-example} 的解是唯一的.
    \end{enumerate}
\end{BookExercises}
